\chapter{系统概述与总体方案}

\section{赛题任务与设计目标}

本项目面向飞跃雷区组空地协同任务。无人机承担环境感知、相对定位、稳定飞行及任务协调，
地面车模完成地面运动与指定目标交互。系统设计目标包括：在限定场地和硬件规则内完成可靠起降，
稳定获取无人机、车模与信标的相对状态，通过低延迟通信形成闭环，并在异常状态下保持可诊断、可急停。

\todoitem{引用正式比赛规则，逐条写清起点、雷区/信标布置、任务顺序、尺寸与芯片约束、计分方法及安全要求。不要仅描述“完成任务”。}

\section{系统总体架构}

当前仓库的主要实现集中于无人机端，其软件架构如图\ref{fig:system-architecture}所示。

\begin{figure}[htbp]
  \centering
  \begin{tikzpicture}[
    node distance=7mm and 9mm,
    box/.style={draw,rounded corners,align=center,minimum height=10mm,minimum width=30mm,fill=blue!4},
    core/.style={box,fill=red!6,draw=BoschRed},
    line/.style={-{Latex[length=2mm]},thick}
  ]
    \node[box] (camera) {双摄像头\\红外/车标图像};
    \node[box,right=of camera] (imu) {ICM42688\\姿态观测};
    \node[box,right=of imu] (range) {双ToF/光流\\高度与速度};

    \node[core,below=of camera] (cm71) {CM7\_1 相机核\\采集与视觉前处理};
    \node[core,below=of imu] (cm70) {CM7\_0 控制核\\估计、控制与任务管理};
    \node[box,below=of range] (carlink) {地面车模链路\\位置、航向与速度};

    \node[box,below=of cm71] (vision) {标志检测、IPM\\双相机融合};
    \node[box,below=of cm70] (control) {位置→速度→姿态\\→角速度级联控制};
    \node[box,below=of carlink] (safety) {安全状态机\\VOFA诊断与遥测};

    \node[box,below=of control,minimum width=42mm] (motor) {姿态优先混控\\DShot600电机输出};

    \draw[line] (camera) -- (cm71);
    \draw[line] (imu) -- (cm70);
    \draw[line] (range) -- (cm70);
    \draw[line] (cm71) -- node[left,font=\small]{跨核数据} (vision);
    \draw[line] (vision) -- (control);
    \draw[line] (cm70) -- (control);
    \draw[line] (carlink) -- (cm70);
    \draw[line] (control) -- (motor);
    \draw[line] (safety) -- (control);
  \end{tikzpicture}
  \caption{基于当前仓库实现整理的系统总体架构}
  \label{fig:system-architecture}
\end{figure}

\section{双核软件分工}

控制核入口为\codepath{project/user/main_cm7_0.c}，负责初始化、菜单与监视、状态管理以及实时控制链路；
相机核入口为\codepath{project/user/main_cm7_1.c}，通过
\codepath{project/code/communication/cross_core_camera.c}承担摄像头任务和跨核传输。
高频控制主要由定时中断驱动，核心调度位于\codepath{project/code/common/PIT_Demo.c}，不能仅以主循环调用顺序描述系统实时性。

\section{项目特色与拟突出创新点}

\begin{enumerate}
  \item \textbf{双核异构分工：}视觉采集与飞行控制分核执行，降低图像任务对高频控制周期的干扰。
  \item \textbf{多源定位链路：}摄像头标定、IPM、红外信标、车标方向识别、光流和ToF共同提供位置、速度与高度信息。
  \item \textbf{分层闭环：}位置环输出目标速度，速度环输出目标加速度/姿态，姿态环和角速度环最终驱动电机。
  \item \textbf{连续轨迹与制动：}位置控制包含P/Sqrt纠偏、速度/加速度/Jerk整形、动态停车距离和快速反向处理。
  \item \textbf{工程诊断与安全：}VOFA多模式观测、控制实时性统计、电池和电调遥测、缓启动及急停锁存形成完整调试闭环。
\end{enumerate}

\todoitem{创新点最终必须由对比试验支撑。为每项创新安排基线、评价指标和消融实验，不使用“零延迟”“完全消除”等绝对表述。}
